\chapter{模型服务能力与 Heptalogos 边界}

\section{按能力合同接入模型}

模型通过明确的能力合同提供：
\begin{itemize}
\item 向量嵌入（Embedding）与重排（Reranking）；
\item 结构化提案生成（StructuredProposal）；
\item 文本、图像和语音派生（Text/Vision/Speech Derivation）。
\end{itemize}
合同描述输入限制、成本、隐私、可用性和 producer identity。向量空间标识包含模型/权重、任务、tokenizer/预处理、pooling、normalization 和 dimension，避免把不同空间的向量混为一体。\cite{nous-research-package}

本地模型网关和远程 API 在 cognition architecture 中承担相同 provider role。部署元数据可以记录时延、隐私边界、离线可用性、输入限制、成本和硬件要求；切换 provider 不改变 SubjectId、Memory identity 或 owner semantics。

\section{Nous Wave 与 Heptalogos 的职责}

目标集成中，Nous Wave 负责主体 cognition/personality：Memory、Self、Social、Motivation、Cognitive Runtime、retrieval 与 offline cognition proposal。Heptalogos 负责 real-IM interaction、canonical product facts、Opportunity/Reaction、行为审查、规范/管理权限、binding Commitment、CommunicationCommit、Action/Effect 和外部世界结果。双方共享逻辑 SubjectId，并通过类型明确的 observation、context/query、proposal、use/outcome receipt 协作。\cite{nous-research-package,heptalogos-code-snapshot}

核心边界为：

\begin{center}
Nous：``我记得/如何理解/我是谁/我与谁是什么关系/我想要什么''\\
Heptalogos：``现在是否行动、是否允许、承诺什么、实际向世界做了什么''
\end{center}

Motivation 中的 Goal 不自动成为 Heptalogos Commitment；Social/Memory 的判断不自动获得行为权限；Heptalogos 的 prohibitions 和 system/admin constraints 也不应被普通 cognition revision 覆盖。

\section{当前 Heptalogos 实现边界}

当前权威实现基线是个人仓库 \texttt{ArsvineZhu/Heptalogos}。该仓库已进入 headless Product Host / Product-Subject path，Foundation 以 Service/Capability、Context/Activity contribution、Evidence/provenance 和 authority ceiling 为高级认知保留接入面；它明确不冻结 Memory、Persona、Relationship 等高级 cognition 的内部 ontology。\cite{heptalogos-code-snapshot}

当前 Heptalogos 架构中的 OpenClaw 属于独立 Machine Operations Plane，用于机器/部署维护能力，不是 Subject cognition runtime，也不是 Nous Wave 的替代实现。Vault 旧材料中把 Subject cognition 与 OpenClaw 直接绑定的描述不再代表当前产品架构。

Nous Wave 与 Heptalogos 尚未完成直接产品集成。后续需要冻结的最小跨库合同包括：

\begin{itemize}
\item SubjectId provisioning/binding；
\item Heptalogos canonical interaction fact $\rightarrow$ Nous Observation；
\item Nous ContextContribution/CognitiveQuery 在 Reaction/Activity 中的消费；
\item Nous Goal/proposal $\rightarrow$ Heptalogos review/adoption；
\item Heptalogos use/outcome $\rightarrow$ Nous 的幂等回执；
\item Heptalogos maintenance opportunity/lease $\rightarrow$ Nous offline cognition workflow。
\end{itemize}

两侧不进行跨数据库分布式 transaction；可靠交接依赖稳定 event/request identity、idempotency、明确的 point-of-no-return 和可重试 receipt。

\section{维护窗口与离线认知}

Nous Wave 不拥有产品级 cron/nightly policy。Heptalogos 或其他 Host 决定“现在是否允许维护”以及可用时间、token/compute、网络和 provider 权限；Nous Wave 根据当前 cognition debt 选择 Journal maintenance、Episode review、Dream/Hypothesis 或其他 owner-specific工作，并在预算耗尽或收益不足时停止。维护窗口提供机会和资源，不赋予 Dream 越权写入 Memory/Self/Social/Motivation 的权限。

\section{权限与模型服务失败}

模型服务输出默认形成 proposal、derived representation 或 retrieval evidence，而不是 Authority。Kernel/owner commit path 重新核验 Subject、evidence lineage、expected revision/epoch、ACL 和生命周期后才能提交。必需 capability 不可用时明确失败；可选 retrieval lane/provider 不可用时可以降级并保留 degradation diagnostics。

PC-local-first 是默认部署目标，但不是“只能本地”。官方 profile 可以使用本地 gateway，也可以使用远程 API；两者通过同一 capability contract 接入。敏感材料是否允许远程调用由 deployment/privacy policy 决定，而不是由 cognition object type 隐式决定。
